\subsection{Definition of a measure}

DEFINITION 2. --- A continuous linear form on \(\mathcal{K}(\mathrm{X};\mathbf{C})\) , \(\mathrm{X}\) a locally compact space, is called a measure (or complex measure) on \(\mathrm{X}\) .

If \(\mu\) is a measure on a locally compact space \(X\) , the value of the measure for a function \(f \in \mathcal{K}(X; \mathbf{C})\) is called the integral of \(f\) with respect to \(\mu\) ; besides the general notations \(\mu(f)\) and \(\langle f, \mu \rangle\) , one also uses the notations \(\int f d\mu, \int f\mu, \int f(x) d\mu(x)\) and \(\int f(x)\mu(x)\) to denote it; as to the use of the letter \(x\) , see S, I, §1, No.~1.

By virtue of the criterion for continuity in direct limits (TVS, Ch. II, §4, No.~4, Prop. 5), to say that \(\mu\) is a measure on X means that \(\mu\) is a linear form on \(\mathcal{K}(\mathrm{X};\mathbf{C})\) satisfying the following condition: for every compact subset K of X, there exists a number \(\mathrm{M_K}\) such that, for every function \(f\in \mathcal{K}(\mathrm{X};\mathbf{C})\) whose support is contained in K,

\[
| \mu (f) | \leqslant \mathrm{M} _ {\mathrm{K}} \cdot \| f \| \quad (\text { where } \| f \| = \sup _ {x \in \mathrm{X}} | f (x) |).\tag{4}
\]

More generally:

PROPOSITION 6. --- Let X be a locally compact space, \((\mathrm{K}_{\alpha})\) a family of compact subsets of X whose interiors \(\mathring{K}_{\alpha}\) form an open covering of X. For a linear form \(\mu\) on \(\mathcal{K}(\mathrm{X};\mathbf{C})\) to be a measure on X, it is necessary and sufficient that, for every \(\alpha\) , there exist a number \(M_{\alpha}\) such that

\[
| \mu (f) | \leqslant \mathrm{M} _ {\alpha} \cdot \| f \|\tag{5}
\]

for every function \(f \in \mathcal{K}(\mathrm{X}, \mathrm{K}_{\alpha}; \mathbf{C})\) .

The condition being obviously necessary, it suffices to prove that (5) implies (4) for every compact subset K of X. Now, K is covered by a finite number of open sets \(\mathring{K}_{\alpha_{i}}\) ( \(1 \leqslant i \leqslant n\) ); applying Lemma 1 of No.~2 to K and to the \(\mathring{K}_{\alpha_{i}}\) , there exist continuous functions \(g_{i} \geqslant 0\) on X such that \(\operatorname{Supp}(g_{i}) \subset \mathrm{K}_{\alpha_{i}}\) , \(0 \leqslant \sum_{i=1}^{n} g_{i}(x) \leqslant 1\) for all \(x \in X\) and \(\sum_{i=1}^{n} g_{i}(x) = 1\) for \(x \in K\) . For every function \(f \in \mathcal{K}(X, K; \mathbf{C})\) , we can therefore write \(f = \sum_{i=1}^{n} fg_{i}\) and we have \(fg_{i} \in \mathcal{K}(X, K_{\alpha_{i}}; \mathbf{C})\) and \(\|fg_{i}\| \leqslant \|f\|\) ; if \(M_{K} = \sum_{i=1}^{n} M_{\alpha_{i}}\) , we then have the relation (4).

We denote by \(\mathcal{M}(\mathrm{X};\mathbf{C})\) , or simply \(\mathcal{M}(\mathrm{X})\) if no confusion can result, the vector space of measures on \(\mathrm{X}\) , in other words, the dual of \(\mathcal{K}(\mathrm{X};\mathbf{C})\) .

One knows that for every set \(\mathfrak{S}\) of bounded subsets of \(\mathcal{K}(\mathrm{X};\mathbf{C})\) , there is defined on \(\mathcal{M}(\mathrm{X};\mathbf{C})\) the \(\mathfrak{S}\) -topology, which is locally convex (TVS, III, §3, No.~1, Cor. of Prop. 1). We denote the topological vector space, obtained by equipping \(\mathcal{M}(\mathrm{X};\mathbf{C})\) with the \(\mathfrak{S}\) -topology, by \(\mathcal{M}_{\mathfrak{S}}(\mathrm{X};\mathbf{C})\) or \(\mathcal{M}_{\mathfrak{S}}(\mathrm{X})\) .

PROPOSITION 7. --- For every set S of bounded subsets of \(\mathcal{K}(\mathrm{X};\mathbf{C})\) that is a covering of \(\mathcal{K}(\mathrm{X};\mathbf{C})\) , the space \(\mathcal{M}_{\mathfrak{S}}(\mathrm{X};\mathbf{C})\) is Hausdorff and quasi-complete.

This results from the fact that \(\mathcal{K}(\mathrm{X};\mathbf{C})\) is barreled (TVS, III, §4, No.~2, Cor. 4 of Th. 1).

Examples of measures. --- I. Atomic measures. Let X be a locally compact space, a a point of X; the mapping \(f \mapsto f(a)\) of \(\mathcal{K}(\mathrm{X};\mathbf{C})\) into C obviously satisfies the condition (4) with \(M_{K} = 1\) for every compact subset K of X containing a, hence is a measure on X, which is denoted by \(\varepsilon_{a}\) ; it is called the Dirac measure at the point a, or the measure defined by a unit mass placed at the point a.

More generally, let \(\alpha\) be a mapping of X into C such that, for every compact subset K of X, \(\sum_{x\in K}|\alpha(x)|<+\infty\) . Then, for every function \(f\in\mathcal{K}(X,K;\mathbf{C})\) , the sum

\[
\mu (f) = \sum_ {x \in X} \alpha (x) f (x)
\]

is defined, being equal to \(\sum_{x\in K}\alpha (x)f(x)\) ; it is clear that \(\mu\) is a linear form on \(\mathcal{K}(\mathrm{X};\mathbf{C})\) and that, for \(f\in \mathcal{K}(\mathrm{X},\mathrm{K};\mathbf{C})\) ,

\[
| \mu (f) | \leqslant \left(\sum_ {x \in K} | \alpha (x) |\right) \cdot \| f \|,
\]

in other words the condition (4) is satisfied.

A measure \(\mu\) on X is said to be atomic if there exists a mapping \(\alpha\) of X into C such that \(\sum_{x\in K}|\alpha (x)| < +\infty\) for every compact subset K of X, and such that \(\mu\) is equal to the measure defined as above. If N is the set of \(x\in X\) such that \(\alpha (x)\neq 0\) , the condition imposed on \(\alpha\) implies that for every compact subset K of X, \(\mathrm{K}\cap \mathrm{N}\) is countable. One also says that \(\mu\) is defined by the masses \(\alpha (x)\) placed at the points \(x\in \mathrm{N}\) . If one assumes that \(\mathrm{N}\cap \mathrm{K}\) is finite for every compact set \(\mathrm{K}\subset \mathrm{X}\) , then obviously \(\sum_{x\in K}|\alpha (x)| < +\infty\) ; it comes to the same to say that N is a closed and discrete subspace of X, because then every point of X has a compact neighborhood containing only a finite number of points of N, and conversely, if this is the case, then every compact subset of X can be covered by a finite number of such neighborhoods. When N is closed and discrete, every atomic measure defined by a function \(\alpha\) such that \(\alpha(x)=0\) on CN is called a discrete measure on X (cf.~§2, No.~5).

\begin{enumerate}
\def\labelenumi{\Roman{enumi}.}
\setcounter{enumi}{1}
\tightlist
\item
  Lebesgue measure. For every function \(f \in \mathcal{K}(\mathbf{R}; \mathbf{C})\) , there exists a compact interval \([a, b]\) of R outside of which f is zero. The integral
\end{enumerate}

\[
\operatorname{I} (f) = \int_ {- \infty} ^ {+ \infty} f (x) d x = \int_ {a} ^ {b} f (x) d x
\]

is therefore defined; moreover, by the theorem of the mean (FRV, II, §1, No.~5, Prop. 6), we have \(|\mathrm{I}(f)| \leqslant (b - a) \| f \|\) ; this shows that \(f \mapsto \mathrm{I}(f)\) is a measure on \(\mathbf{R}\) , which is called Lebesgue measure.

For every interval J (bounded or not) of R, one similarly calls Lebesgue measure on J the measure \(f \mapsto \int_{J} f(x) dx\) , a linear form on \(\mathcal{K}(J; \mathbf{C})\) (the integral having meaning since there exists a compact interval \([a, b]\) contained in J outside of which f is zero).

\begin{enumerate}
\def\labelenumi{\Roman{enumi}.}
\setcounter{enumi}{2}
\tightlist
\item
  Let g be a continuous mapping of a compact interval \(I \subset R\) into C, having a continuous derivative in I. Let \(\Gamma = g(I)\) , which is a compact subspace of C; the mapping
\end{enumerate}

\[
f \mapsto \int_ {\mathrm{I}} f (g (t)) g ^ {\prime} (t) d t
\]

of \(\mathcal{C}(\Gamma;\mathbf{C})\) into C is a continuous linear form by virtue of the theorem of the mean, hence is a complex measure on \(\Gamma\) ; the integral relative to this measure is also written \(\int_{\Gamma} f(z) dz\) , even though it depends not only on \(\Gamma\) but also on g.

Remark. --- The giving of a measure \(\mu\) on a locally compact space X defines on X (along with the topology of X) a structure S. Let \(X_{1}\) be a second set, \(\varphi\) a bijective mapping of X onto \(X_{1}\) ; in conformity with general definitions (S, R, §8), the structure \(S_{1}\) obtained by transporting to \(X_{1}\) the structure S of X, by means of \(\varphi\) , is defined in the following way. The topology of X is transported to \(X_{1}\) by \(\varphi\) ; the functions of \(\mathcal{K}(X_{1};\mathbf{C})\) are then the functions f such that \(f \circ \varphi\) belongs to \(\mathcal{K}(X;\mathbf{C})\) , and the measure \(\mu_{1}\) on \(X_{1}\) is defined by \(\mu_{1}(f) = \mu(f \circ \varphi)\) .

In particular, an automorphism of the structure S is a homeomorphism \(\sigma\) of X onto itself, such that

\[
\mu (f) = \mu (f \circ \sigma)
\]

for every function \(f \in \mathcal{K}(\mathrm{X};\mathbf{C})\) ; the measure \(\mu\) is then also said to be invariant under the homeomorphism \(\sigma\) .

Example. --- Lebesgue measure on \(\mathbf{R}\) is invariant under every translation of the additive group \(\mathbf{R}\) . Indeed, for every function \(f \in \mathcal{K}(\mathbf{R}; \mathbf{C})\) and every real number \(a\) , we have

\[
\int_ {- \infty} ^ {+ \infty} f (x + a) d x = \int_ {- \infty} ^ {+ \infty} f (t) d t
\]

by the change-of-variables formula (FRV, II, §2, No.~1, formula (1)).

For a generalization, see Ch. VII, §1, No.~2, Th. 1.

